\subsection{Almost abelian solvmanifolds}\label{sec:abeliansolv}

In contrast to the case of nilpotent groups, there is no
simple criterion for the existence of a~lattice in a~general
connected and simply-connected solvable Lie group, as is required to define a corresponding twisted
torus. To formulate such a criterion, we specialise to \emph{almost
abelian solvable groups}: these are the solvable Lie groups $\sfG$ of
dimension $d$
whose nilradical $\sfN$ has codimension one and is abelian:
\[
\sfN \simeq \real^{d-1} .
\]
Then $\sfG$ has the
structure of a semi-direct product
\[
\sfG = \sfN\rtimes_\varphi\real
\]
for a continuous one-parameter left group action
$\varphi\colon \real\to\Aut(\sfN)$; concretely, $\varphi$ is given by the
adjoint action of the one-dimensional subgroup $\sfH=\real$ on $\sfN$
in the group $\sfG$ (cf.~Section~\ref{sec:semidirect}). This exhibits
$\sfG$ as a nontrivial group extension
\[
1\longrightarrow \sfN \longrightarrow \sfG \longrightarrow \R \longrightarrow 1 .
\]

We can regard the one-parameter group action $\varphi$ as a matrix
$\varphi_x\in\sfGL(d-1,\real)$ for each $x\in\real$ acting on the
vector space $\sfN$, which we identify with $\real^{d-1}$ via a choice
of basis $\vec e_1,\ldots,\vec e_{d-1}$. Since
$\varphi_{0}=\unit_{d-1}$ and $\varphi_x$ is always non-singular, it
follows from the continuity of $\varphi$ and the determinant that
$\det\varphi_x>0$ for all $x\in\real$. Then $\sfG$ admits a lattice
$\Lambda_\sfG$ if and
only if there exists $x_0\in\real^\times$ such that $\varphi_{x_0}$ is
conjugate to an integer matrix $\ttM\in\sfSL(d-1,\zed)$:
\begin{gather}\label{eq:monodromydef}
\Sigma^{-1} \varphi_{x_0} \Sigma=\ttM
\end{gather}
for some
$\Sigma\in\sfGL(d-1,\real)$.
In this case the twisted torus
$\torus_{\Lambda_\sfG}=\Lambda_{\sfG}\backslash \sfG$ is called an {\emph{almost abelian solvmanifold}}.
The condition~\eqref{eq:monodromydef} strongly restricts the homomorphisms
$\varphi\colon \real\to\Aut(\sfN)$; in particular, it requires that the
characteristic polynomial of $\varphi_{x_0}$ has integer
coefficients. In this case the lattice (in the standard basis $\vec
e_1,\ldots,\vec e_{d-1}$) is given by
\[%\label{standardbasis}
\Lambda_\sfG = \Sigma\cdot\zed^{d-1}\rtimes_{\varphi|_{x_0\,\zed}}
x_0 \zed ,
\]
which correspondingly sits as a nontrivial group extension
\[
1\longrightarrow \Sigma \cdot \zed^{d-1} \longrightarrow \Lambda_\sfG
\longrightarrow x_0 \zed \longrightarrow 1 .
\]
Then $\Lambda_{\sfN}\setminus\sfN\simeq\torus^{d-1}$, and the corresponding Mostow bundle realises the twisted torus
$\torus_{\Lambda_\sfG}$ as a torus bundle over a circle
$\Lambda_\sfG \sfN\backslash\sfG \simeq \torus$, whose monodromy
is specified by the matrix $\ttM$ in the mapping class group $\sfSL(d-1,\zed)$ of
orientation-preserving
automorphisms up to homotopy of the torus fibres~$\torus^{d-1}$.

For an almost abelian solvmanifold we can make the twisted torus construction more concrete
by choosing the global
coordinates $(z,x)\in\real^{d-1}\times\real$ on the group
manifold (associated with the basis $\vec e_1,\ldots,\vec e_{d-1}$).
The group multiplication of the semi-direct product
$\sfG=\real^{d-1}\rtimes_\varphi\real$ is
\begin{gather}\label{eq:grouplaw}
(z,x) (z',x') = (z + \varphi_x\cdot z',x+x') ,
\end{gather}
where we used $\varphi_x \varphi_{x'} = \varphi_{x+x'}$, and the
inverse of a group element is
\[
(z,x)^{-1} = (-\varphi_{-x}\cdot z,-x) ,
\]
where we used $\varphi_x^{-1}=\varphi_{-x}$.
The twisted torus is defined as the quotient $\torus_{\Lambda_\sfG}=\Lambda_\sfG
{\setminus} \sfG$ which is generated by the equivalence relation $(z,x)\sim
(\Sigma\cdot\gamma,x_0\,\alpha)\,(z,x)$ for all $(\gamma,\alpha)\in\zed^{d-1}\times\zed$. The global structure
of $\torus_{\Lambda_{\sfG}}$ is generated by the simultaneous local coordinate
identifications under the action of the elements
$(\Sigma\cdot\gamma,x_0 \alpha) $ of $\Lambda_\sfG$ given by
\begin{gather}
(z,x) \longmapsto (
 z+\Sigma\cdot\gamma,x) , \nonumber \\
(z,x) \longmapsto (
\Sigma {\tt M}^\alpha\,\Sigma^{-1}\cdot z,x+x_0\,\alpha) .\label{eq:Mostowtorus}
\end{gather}
These identifications explicitly exhibit the twisted torus as a torus bundle over a circle, with
local fiber coordinates $z\in \torus^{d-1}$ and base coordinate $x\in \torus$,
 whose monodromy is specified by the matrix ${\tt M}\in
 \sfSL(d-1,\zed)$, and whose periods are given respectively by
 $\Sigma\in\sfGL(d-1,\real)$ and $x_0\in\real^\times$.

\subsection{$C^*$-algebra bundles}\label{sec:C*bundles}

Mostow bundles and their $C^*$-algebraic T-duals can be grouped together
under the general heading of `$C^*$-algebra bundles', which encompasses
the notions of $C_0(X)$-algebras and $C^*$-bundles, as we now
explain; see Section~8.1 and Appendix~C of~\cite{Williams2007} for
further details.
Let~$X$ be a locally compact Hausdorff space. There are two equivalent
notions for the $C^*$-algebra analogue of a fibre bundle over $X$.

A \emph{$C_0(X)$-algebra} is a $C^*$-algebra $\alg$
equipped with a nondegenerate injection $\iota$ of $C_0(X)$
into the centre of its multiplier algebra, called the \emph{structure
 map}. For $\mathtt{f}\in C_0(X)$ and
$a\in\alg$, we abbreviate $\iota(\mathtt{f})a$ by $\mathtt{f}\cdot a$. This endows
$\alg$ with a $C_0(X)$-bimodule
structure.

Given a family $\balg=(\balg_x)_{x\in X}$ of $C^*$-algebras, a
\emph{section} of $\balg$ is
a map $s\colon X\to\balg$ such that $s(x)\in\balg_x$ for all $x\in X$; we
denote the space of sections of $\balg$ which vanish at infinity by~$\Gamma_0(\balg)$. The family $\balg$ is then called a~\emph{$C^*$-bundle} over $X$ with fibres $\balg_x$ if the following
conditions are satisfied:
\begin{itemize}\itemsep=0pt
\item
$\Gamma_0(\balg)$ is a $C^*$-algebra under pointwise operations and the
 supremum norm;
\item $\balg_x=\{s(x)\, |\, s\in\Gamma_0(\balg)\}$ for each
 $x\in X$;
\item $\Gamma_0(\balg)$ is closed under multiplication by
 $C_0(X)$; and
\item For each $s\in\Gamma_0(\balg)$, the function $x\mapsto \|s(x)\|$
is upper semi-continuous, i.e., the set
$\{x\in X \,|\, \|s(x)\|< \epsilon\}$ is open in $X$ for all $\epsilon>0$.
\end{itemize}
In this
paper we will only be concerned with $C^*$-bundles that have non-zero fibres.

If $\balg$ is a $C^*$-bundle over $X$, then its section algebra
$\Gamma_0(\balg)$ is a $C_0(X)$-algebra: its structure map $\iota$ from
$C_0(X)$ is defined by $\iota(\mathtt{f})s=\mathtt{f} s$. Conversely, if $\alg$ is a
$C_0(X)$-algebra, then the fibre~$\alg_x$ of $\alg$ over $x\in X$ is $\alg_x:=\alg/\CI_x$, where
$\CI_x=\{\mathtt{f}\cdot a\, |\, \mathtt{f}\in C_0(X) , \,
\mathtt{f}(x)=0 , \, a\in\alg\}$ is identified as the ideal in
$\alg$ of sections vanishing at $x$. If
$a\in\alg$, we write $a(x)=a+\CI_x$ for its image in~$\alg_x$. The
function $x\mapsto\|a(x)\|$ is upper semi-continuous
and vanishes at infinity with
\[
\|a\| = \sup_{x\in X} \big\|a(x)\big\|
\]
for all $a\in\alg$. The
elements $a\in \alg$ can in this way be viewed as sections of a
$C^*$-bundle $(\alg_x)_{x\in X}$. We will sometimes use the
notation $\coprod_{x\in X} \alg_x$ for $\alg$ when we wish to emphasise
its structure as a $C^*$-bundle over $X$ with fibre $C^*$-algebras
$\alg_x$. These definitions do not require
local triviality of the bundle nor the fibres of the bundle to be
isomorphic to one another. $C^*$-algebra bundles over $X$ are objects of a
category whose morphisms are \emph{fibrewise} $*$-homomorphisms,
i.e., $C_0(X)$-linear morphisms $\psi\colon \alg\to\balg$: $\psi(\mathtt{f}\cdot
a)=\mathtt{f}\cdot\psi(a)$ for all $\mathtt{f}\in C_0(X)$ and $a\in\alg$; then $\psi$
induces $*$-homomorphisms $\psi_x\colon \alg_x\to\balg_x$ such that
$\psi_x\big(a(x)\big)=\psi(a)(x)$ for all $a\in\alg$.

\begin{Example}[trivial $C^*$-algebra bundles]
If $\Dcal$ is any $C^*$-algebra, then $\alg=C_0(X,\Dcal)\simeq
C_0(X)\otimes\Dcal$ is naturally a $C_0(X)$-algebra with structure map
\[
\big(\iota(\mathtt{f})a\big)(x) := \mathtt{f}(x)\,a(x)
\]
for $\mathtt{f}\in C_0(X)$, $a\in\alg$ and $x\in X$. In this case each fibre
$\alg_x$ is canonically identified with~$\Dcal$ and elements of~$\alg$
are obviously identified with sections.
\end{Example}

\begin{Example}[continuous maps]\label{ex:ctsmaps}
Let~$X$ and~$Y$ be locally compact spaces and $\sigma\colon Y\to X$ a~continuous surjective map. Then $C_0(Y)$ is a $C_0(X)$-algebra with structure map
$\iota(\mathtt{f})\mathtt{g}:=(\mathtt{f}\circ\sigma)\,\mathtt{g}$, for $\mathtt{f}\in C_0(X)$ and $\mathtt{g}\in C_0(Y)$,
and fibers $C_0(Y)_x\simeq C_0\big(\sigma^{-1}(x)\big)$.
\end{Example}

In this paper we are particularly interested in crossed
products of $C^*$-algebra bundles. Let~$\alg$ be a $C_0(X)$-algebra,
and denote by $\Aut_X(\alg)$ the group of {fibrewise
automorphisms} of~$\alg$. A~\emph{fibrewise action} of a locally compact group $\sfG$ on
$\alg$ is then a group homomorphism $\alpha\colon \sfG\to\Aut_X(\alg)$. This
implies that $\alpha$ induces an action~$\alpha^x$ on each fiber
$\alg_x$ for $x\in X$, and in
this case we say that the dynamical system $(\alg,\sfG,\alpha)$ is
\emph{$C_0(X)$-linear}.

\begin{Theorem}\label{thm:crossedfibres}
Let $X$ be a locally compact Hausdorff space, and let
$(\alg,\sfG,\alpha)$ be a $C_0(X)$-linear $C^*$-dynamical system. Then
the crossed product $\alg\rtimes_\alpha\sfG$ is again a
$C_0(X)$-algebra with fibres
\[
(\alg\rtimes_\alpha\sfG)_x\simeq\alg_x\rtimes_{\alpha^x}\sfG ,
\]
where
\[
\alpha_\gamma^x\big(a(x)\big) = \alpha_\gamma(a)(x)
\]
for each $x\in X$, $\gamma\in\sfG$ and $a\in\alg$.
\end{Theorem}

\begin{proof}The structure map
of $\alg\rtimes_\alpha\sfG$ is given by precomposing the structure map
of $\alg$ with the natural injection of the center of the multiplier
algebra of $\alg$ into the center of the multiplier algebra of
$\alg\rtimes_\alpha\sfG$; it satisfies
$(\mathtt{f}\cdot f)(\gamma)=\mathtt{f}\cdot\big(f(\gamma)\big)$ for all
$\mathtt{f}\in C_0(X)$, $f\in C_{\rm c}(\sfG,\alg)$ and
$\gamma\in\sfG$. See~\cite[Theorem~8.4]{Williams2007} for further
details.
\end{proof}

